A $UBV$ photometric ($UBV$ Johnson's) observation of a star gives
$U = 8.15$, $B = 8.50$, and $V = 8.14$. Based on the spectral class, one gets
the intrinsic color $(U-B)_{\mathrm{o}} = -0.45$. If the star is known to
have radius of $2.3\,R_\odot$, absolute bolometric magnitude of $-0.25$, and
bolometric correction ($BC$) of $-0.15$, determine:
\begin{description}
\item[a.] the intrinsic magnitudes $U$, $B$, and $V$ of the star (take, for
  the typical interstellar matters, the ratio of total to selective
  extinction (color excess) $R_V = 3.2$),
\item[b.] the effective temperature of the star,
\item[c.] the distance to the star in pc.
\end{description}
Note: The relation between color excess of $U-B$ and of $B-V$ is
$E(U-B) = 0.72\,E(B-V)$.
